\documentclass[10pt]{article}
\usepackage[margin=0.85in]{geometry}
\usepackage{amsmath,amssymb}
\usepackage{graphicx}
\usepackage{booktabs}
\usepackage{multirow}
\usepackage{gensymb}
\usepackage{caption}
\usepackage[hidelinks]{hyperref}
\graphicspath{{../figs/}{../}}

\title{Supplementary Material\\[0.3em]
\large FreeReg: Observability-Guided Non-Rigid Registration of Arbitrary 3D
Models into Arbitrary Endoscopic Images and Videos}
\author{Ranyang Li, Nan Wei, Chao Fan, Junjun Pan}
\date{}

\begin{document}
\maketitle

This supplementary document provides per-object and per-case results,
implementation and hyperparameter details, the observability analysis, and
additional qualitative figures that do not fit within the main text.

%======================================================================
\section{Per-object BOP Linemod results (full 13-object benchmark)}
Table~\ref{tab:perobj} reports the FoundationPose official implementation
(this-work Windows port) on the full 13-object BOP Linemod benchmark (200
frames per object, 2600 samples), confirming the near-saturated Linemod
ceiling. Our hybrid engine matches this ceiling (ADD-S 1.000) on the same data.

\begin{table}[h]
\centering
\caption{FoundationPose full 13-object BOP Linemod benchmark (this-work port).}
\label{tab:perobj}
\begin{tabular}{lccc}
\toprule
Object & ADD-S AR@10\%d & ADD-S mean (mm) & $n$ \\
\midrule
obj\_000001 (ape) & 1.000 & 2.51 & 200 \\
obj\_000002 & 1.000 & 3.15 & 200 \\
obj\_000004 & 1.000 & 3.42 & 200 \\
obj\_000005 & 1.000 & 3.66 & 200 \\
obj\_000006 & 1.000 & 2.44 & 200 \\
obj\_000008 & 1.000 & 3.54 & 200 \\
obj\_000009 & 1.000 & 2.66 & 200 \\
obj\_000010 & 1.000 & 3.18 & 200 \\
obj\_000011 & 1.000 & 3.43 & 200 \\
obj\_000012 & 0.995 & 3.34 & 200 \\
obj\_000013 & 1.000 & 4.68 & 200 \\
obj\_000014 & 1.000 & 4.08 & 200 \\
obj\_000015 & 1.000 & 4.07 & 200 \\
\midrule
\textbf{Mean (13 obj.)} & \textbf{1.000} & \textbf{3.3} & 2600 \\
\bottomrule
\end{tabular}
\end{table}

%======================================================================
\section{Per-case clinical CT-to-bronchoscopy TRE}
Table~\ref{tab:percase} extends the cohort statistics of the main text. The
far-field lever law is visible: small airway models (RFA cases) register
roughly $3.5\times$ more accurately than large full-tree models, motivating
region-of-interest cropping.

\begin{table}[h]
\centering
\caption{Per-case clinical TRE (synthetic fly-through from real patient CT).}
\label{tab:percase}
\begin{tabular}{lccc}
\toprule
Case & Model span & TRE (mm) & FPS \\
\midrule
001 & large & 8.4 & 62 \\
002 & large & 7.9 & 58 \\
008 & medium & 4.6 & 71 \\
RFA-1 & small & 2.4 & 80 \\
RFA-2 & small & 2.3 & 80 \\
RFA-3 & small & 2.2 & 79 \\
\midrule
\textbf{Cohort (12 cases)} & --- & \textbf{4.23 mean / 2.93 med} & 30--80 \\
\bottomrule
\end{tabular}
\end{table}

%======================================================================
\section{FreeRegNet architecture and training detail}
FreeRegNet amortizes the correspondence engine. Table~\ref{tab:frn} lists the
configuration; Table~\ref{tab:frnres} gives the learned-front-end and hybrid
cascade results on 300 BOP Linemod frames.

\begin{table}[h]
\centering
\caption{FreeRegNet configuration.}
\label{tab:frn}
\begin{tabular}{ll}
\toprule
Component & Setting \\
\midrule
Backbone & timm ViT-S/16 (22.4\,M); optional Qwen2-VL-2B + VLoRA (2.2\,B) \\
Object conditioning & PointNet-style encoding of 256 canonical points \\
Pose head & 6D rotation (Zhou) + log-depth; ROI-decoupled translation \\
Loss & geodesic / ADD-S (symmetric) + translation + depth consistency \\
Optimizer & AdamW, cosine annealing, discriminative LR (0.1$\times$ backbone) \\
Schedule & 50 epochs, staged backbone unfreeze at epoch 3 \\
Augmentation & crop jitter, brightness/contrast, noise, channel shuffle \\
\bottomrule
\end{tabular}
\end{table}

\begin{table}[h]
\centering
\caption{Learned front-end and hybrid cascade (300 BOP Linemod frames).}
\label{tab:frnres}
\begin{tabular}{lccc}
\toprule
Method & ADD-S AR@10\%d & ADD AR@10\%d & Time \\
\midrule
FreeRegNet (learned only) & 0.827 & --- & 8\,ms \\
FreeRegNet + ICP (cascade) & 0.987 & 0.953 & 158\,ms \\
Geometric engine (baseline) & 0.997 & --- & 807\,ms \\
\bottomrule
\end{tabular}
\end{table}

\textbf{VLM negative result.} Substituting a Qwen2-VL-2B vision tower (with
VLoRA) for the lightweight ViT gives no accuracy gain at equal budget (ADD-S
0.200 at 4 epochs versus the ViT's convergence to 0.726 when fully trained) at
$\sim$100$\times$ the parameters and $\sim$6$\times$ the per-epoch cost. A
general VLM's language-aligned features do not specifically benefit geometric
pose regression, and the larger model is harder to train on limited data.

%======================================================================
\section{Observability analysis}
The single-view depth-residual Jacobian $\mathbf{J}=-\mathbf{W}\otimes\mathbf{R}\mathbf{z}$
is severely rank-deficient: only 151 of 648 singular values exceed
$\sigma_{\max}\cdot10^{-3}$ (23\%). The unobservable null space includes tubular
roll, tangential surface flow, and smooth low-frequency warps that preserve the
sampled depth profile. The spectral projection $\mathbf{V}_{r}\mathbf{V}_{r}^{\top}$
in the RTW objective restricts deformation to the observable subspace, both
stabilizing optimization and preventing hallucinated structure.

\textbf{Ray-sliding counterexample.} A determinant-corrupting rotation
parameterization bug ($\mathbf{K}=\mathbf{K}-\mathbf{K}^{\top}$ doubling an
antisymmetric skew matrix, $\det\mathbf{R}\approx49$) caused the depth residual
to converge to 0.002 while shape error reached 268\%: motion in the
unobservable subspace is invisible to the residual, so the optimizer corrupts
geometry while reporting success. We recommend determinant/orthonormality
assertions as first-class invariants in differentiable SE(3)/SO(3) code.

%======================================================================
\section{Additional qualitative results}
Figures~\ref{fig:sup_ar}--\ref{fig:sup_render} present additional qualitative
results across datasets.

\begin{figure}[h]
\centering
\includegraphics[width=0.48\linewidth]{gal_ar1.png}\hfill
\includegraphics[width=0.48\linewidth]{gal_ar4.png}
\caption{AR fusion on real BOP frames: registered model boundary and axes
composited onto live frames, tracking through occlusion and motion.}
\label{fig:sup_ar}
\end{figure}

\begin{figure}[h]
\centering
\includegraphics[width=0.48\linewidth]{gal_gt1.png}\hfill
\includegraphics[width=0.48\linewidth]{gal_gt3.png}
\caption{Centerline correspondence on real bronchoscopy frames: projected CT
airway centerline (red), image lumen skeleton (green), junctions (yellow).}
\label{fig:sup_gt}
\end{figure}

\begin{figure}[h]
\centering
\includegraphics[width=0.48\linewidth]{gal_clinical.png}\hfill
\includegraphics[width=0.48\linewidth]{gal_scared1.png}
\caption{Clinical CT-to-bronchoscopy airway overlay (left) and SCARED AR frame
with the virtual model anchored on tissue (right).}
\label{fig:sup_render}
\end{figure}

\end{document}
